% 依据已确认模型设计重建的真实矢量：同一预测、两项处置、不同后果。
\begin{tikzpicture}[figure picture,foundation picture,
 font=\figurefont\mdseries\fontsize{9}{11}\selectfont,
 trust module/.style={draw=FigureStroke,line width=1pt,rounded corners=1.5mm,fill=white},
 trust wire/.style={draw=FigureStroke,line width=.65pt},
 trust icon/.style={draw=FigureStroke,line width=.8pt,line cap=round,line join=round}]
 % 简化对象图标的真实路径；调用方定义trust icon和trust wire。
\providecommand{\TrustVectorMail}[3]{%
 \begin{scope}[shift={(#1,#2)},scale=#3]
  \draw[trust icon,fill=white,rounded corners=.4mm] (0,0) rectangle (8,5.5);
  \draw[trust icon] (0,5.5)--(4,2.3)--(8,5.5) (0,0)--(2.7,2.4) (8,0)--(5.3,2.4);
 \end{scope}}
\providecommand{\TrustVectorFolder}[3]{%
 \begin{scope}[shift={(#1,#2)},scale=#3]
  \draw[trust icon,fill=FigureBlueFill,rounded corners=.4mm] (0,0)--(0,5)--(3,5)--(4,4)--(8,4)--(8,0)--cycle;
 \end{scope}}
\providecommand{\TrustVectorRecord}[3]{%
 \begin{scope}[shift={(#1,#2)},scale=#3]
  \draw[trust icon,fill=white] (0,0)--(6,0)--(6,6)--(4,8)--(0,8)--cycle;
  \draw[trust icon] (4,8)--(4,6)--(6,6);
  \draw[trust icon] (1,4.8)--(4.8,4.8) (1,3.1)--(4.8,3.1) (1,1.4)--(3.8,1.4);
 \end{scope}}
\providecommand{\TrustVectorDocs}[3]{%
 \begin{scope}[shift={(#1,#2)},scale=#3]
  \foreach \dx/\dy in {0/2,1.6/1,3.2/0}{
   \draw[trust icon,fill=white,rounded corners=.5mm] (\dx,\dy) rectangle ({\dx+6},{\dy+8});
  }
  \draw[trust icon] (4.2,6)--(8.2,6) (4.2,4.1)--(8.2,4.1) (4.2,2.2)--(7.2,2.2);
 \end{scope}}
\providecommand{\TrustVectorTrash}[3]{%
 \begin{scope}[shift={(#1,#2)},scale=#3]
  \draw[trust icon,fill=FigureRedFill] (.7,0)--(0,6)--(6,6)--(5.3,0)--cycle;
  \draw[trust icon,fill=white,rounded corners=.3mm] (-.4,6) rectangle (6.4,7);
  \draw[trust icon] (2,7)--(2,8)--(4,8)--(4,7) (1.8,1)--(1.4,5) (3,1)--(3,5) (4.2,1)--(4.6,5);
 \end{scope}}
\providecommand{\TrustVectorMesh}[3]{%
 \begin{scope}[shift={(#1,#2)},scale=#3]
  \foreach \xx/\yy/\id in {-5/-2/a,-5/2/b,0/-4/c,0/0/d,0/4/e,5/-2/f,5/2/g}{\coordinate (icon-\id) at (\xx,\yy);}
  \foreach \i in {a,b}{\foreach \j in {c,d,e}{\draw[trust wire] (icon-\i)--(icon-\j);}}
  \foreach \i in {c,d,e}{\foreach \j in {f,g}{\draw[trust wire] (icon-\i)--(icon-\j);}}
  % 层内同色、层间蓝红黄；连接保持灰色，不把整网染成单色。
  \foreach \id/\col in {a/FigureBlue,b/FigureBlue,c/FigureRed,d/FigureRed,e/FigureRed,f/FigureYellow,g/FigureYellow}{\draw[trust icon,fill=\col] (icon-\id) circle (.7);}
 \end{scope}}

 \path[use as bounding box] (0,0) rectangle (110,56);
 \draw[trust module,fill=white] (0,21) rectangle (18,38);
 \TrustVectorMail{5}{29}{1}
 \node at (9,24) {正常邮件$\bm x$};
 \draw[trust module] (26,14) rectangle (49,45);
 \node at (37.5,40) {分类器$f$};
 \node at (37.5,34) {$p=f(\bm x)$};
 \TrustVectorMesh{37.5}{23}{1.4}
 \draw[trust module,fill=FigureBlueFill] (62,34) rectangle (82,54);
 \node at (72,49) {$d_1$：隔离};
 \TrustVectorMail{68}{40}{.9}
 \TrustVectorFolder{67}{37}{1.1}
 \draw[trust module,fill=FigureRedFill] (62,4) rectangle (82,24);
 \node at (72,19) {$d_2$：删除};
 \TrustVectorTrash{65}{7}{.8}
 \draw[trust icon,fill=white] (74,8)--(77,7)--(79,10)--(76,11)--cycle;
 \draw[trust icon] (74,8)--(77,10)--(79,10);
 \draw[trust module] (88,34) rectangle (110,54);
 \node at (99,49) {原件可找回};
 \node[anchor=west,inner sep=0pt] at (90,38) {收件人};
 \TrustVectorMail{102}{36}{.75}
 \draw[foundation flow,line width=.8pt,-{Stealth[length=1.2mm,width=.8mm]}] (109,41) to[out=0,in=0] (106,44);
 \draw[trust module] (88,4) rectangle (110,24);
 \node at (99,19) {原件已丢失};
 \node[anchor=west,inner sep=0pt] at (90,8) {收件人};
 \draw[trust icon] (105,8) circle (2.5);
 \draw[trust icon] (105,9.8)--(105,8)--(106.5,7);
 \draw[draw=FigureRed,line width=.9pt] (107,4.5)--(109,6.5) (107,6.5)--(109,4.5);
 \draw[foundation flow] (18,29.5)--(26,29.5);
 \draw[draw=FigureStroke,line width=1.1pt] (49,29.5)--(54,29.5);
 \fill[FigureStroke] (54,29.5) circle (.8);
 \draw[foundation flow] (54,29.5) to[out=70,in=180] (62,44);
 \draw[foundation flow] (54,29.5) to[out=-70,in=180] (62,14);
 \draw[foundation flow] (82,44)--(88,44);
 \draw[foundation flow] (82,14)--(88,14);
\end{tikzpicture}%
